% Options for packages loaded elsewhere
\PassOptionsToPackage{unicode}{hyperref}
\PassOptionsToPackage{hyphens}{url}
%
\documentclass[
]{article}
\usepackage{amsmath,amssymb}
\usepackage{iftex}
\ifPDFTeX
  \usepackage[T1]{fontenc}
  \usepackage[utf8]{inputenc}
  \usepackage{textcomp} % provide euro and other symbols
\else % if luatex or xetex
  \usepackage{unicode-math} % this also loads fontspec
  \defaultfontfeatures{Scale=MatchLowercase}
  \defaultfontfeatures[\rmfamily]{Ligatures=TeX,Scale=1}
\fi
\usepackage{lmodern}
\ifPDFTeX\else
  % xetex/luatex font selection
\fi
% Use upquote if available, for straight quotes in verbatim environments
\IfFileExists{upquote.sty}{\usepackage{upquote}}{}
\IfFileExists{microtype.sty}{% use microtype if available
  \usepackage[]{microtype}
  \UseMicrotypeSet[protrusion]{basicmath} % disable protrusion for tt fonts
}{}
\makeatletter
\@ifundefined{KOMAClassName}{% if non-KOMA class
  \IfFileExists{parskip.sty}{%
    \usepackage{parskip}
  }{% else
    \setlength{\parindent}{0pt}
    \setlength{\parskip}{6pt plus 2pt minus 1pt}}
}{% if KOMA class
  \KOMAoptions{parskip=half}}
\makeatother
\usepackage{xcolor}
\usepackage[margin=1in]{geometry}
\usepackage{color}
\usepackage{fancyvrb}
\newcommand{\VerbBar}{|}
\newcommand{\VERB}{\Verb[commandchars=\\\{\}]}
\DefineVerbatimEnvironment{Highlighting}{Verbatim}{commandchars=\\\{\}}
% Add ',fontsize=\small' for more characters per line
\usepackage{framed}
\definecolor{shadecolor}{RGB}{248,248,248}
\newenvironment{Shaded}{\begin{snugshade}}{\end{snugshade}}
\newcommand{\AlertTok}[1]{\textcolor[rgb]{0.94,0.16,0.16}{#1}}
\newcommand{\AnnotationTok}[1]{\textcolor[rgb]{0.56,0.35,0.01}{\textbf{\textit{#1}}}}
\newcommand{\AttributeTok}[1]{\textcolor[rgb]{0.13,0.29,0.53}{#1}}
\newcommand{\BaseNTok}[1]{\textcolor[rgb]{0.00,0.00,0.81}{#1}}
\newcommand{\BuiltInTok}[1]{#1}
\newcommand{\CharTok}[1]{\textcolor[rgb]{0.31,0.60,0.02}{#1}}
\newcommand{\CommentTok}[1]{\textcolor[rgb]{0.56,0.35,0.01}{\textit{#1}}}
\newcommand{\CommentVarTok}[1]{\textcolor[rgb]{0.56,0.35,0.01}{\textbf{\textit{#1}}}}
\newcommand{\ConstantTok}[1]{\textcolor[rgb]{0.56,0.35,0.01}{#1}}
\newcommand{\ControlFlowTok}[1]{\textcolor[rgb]{0.13,0.29,0.53}{\textbf{#1}}}
\newcommand{\DataTypeTok}[1]{\textcolor[rgb]{0.13,0.29,0.53}{#1}}
\newcommand{\DecValTok}[1]{\textcolor[rgb]{0.00,0.00,0.81}{#1}}
\newcommand{\DocumentationTok}[1]{\textcolor[rgb]{0.56,0.35,0.01}{\textbf{\textit{#1}}}}
\newcommand{\ErrorTok}[1]{\textcolor[rgb]{0.64,0.00,0.00}{\textbf{#1}}}
\newcommand{\ExtensionTok}[1]{#1}
\newcommand{\FloatTok}[1]{\textcolor[rgb]{0.00,0.00,0.81}{#1}}
\newcommand{\FunctionTok}[1]{\textcolor[rgb]{0.13,0.29,0.53}{\textbf{#1}}}
\newcommand{\ImportTok}[1]{#1}
\newcommand{\InformationTok}[1]{\textcolor[rgb]{0.56,0.35,0.01}{\textbf{\textit{#1}}}}
\newcommand{\KeywordTok}[1]{\textcolor[rgb]{0.13,0.29,0.53}{\textbf{#1}}}
\newcommand{\NormalTok}[1]{#1}
\newcommand{\OperatorTok}[1]{\textcolor[rgb]{0.81,0.36,0.00}{\textbf{#1}}}
\newcommand{\OtherTok}[1]{\textcolor[rgb]{0.56,0.35,0.01}{#1}}
\newcommand{\PreprocessorTok}[1]{\textcolor[rgb]{0.56,0.35,0.01}{\textit{#1}}}
\newcommand{\RegionMarkerTok}[1]{#1}
\newcommand{\SpecialCharTok}[1]{\textcolor[rgb]{0.81,0.36,0.00}{\textbf{#1}}}
\newcommand{\SpecialStringTok}[1]{\textcolor[rgb]{0.31,0.60,0.02}{#1}}
\newcommand{\StringTok}[1]{\textcolor[rgb]{0.31,0.60,0.02}{#1}}
\newcommand{\VariableTok}[1]{\textcolor[rgb]{0.00,0.00,0.00}{#1}}
\newcommand{\VerbatimStringTok}[1]{\textcolor[rgb]{0.31,0.60,0.02}{#1}}
\newcommand{\WarningTok}[1]{\textcolor[rgb]{0.56,0.35,0.01}{\textbf{\textit{#1}}}}
\usepackage{longtable,booktabs,array}
\usepackage{calc} % for calculating minipage widths
% Correct order of tables after \paragraph or \subparagraph
\usepackage{etoolbox}
\makeatletter
\patchcmd\longtable{\par}{\if@noskipsec\mbox{}\fi\par}{}{}
\makeatother
% Allow footnotes in longtable head/foot
\IfFileExists{footnotehyper.sty}{\usepackage{footnotehyper}}{\usepackage{footnote}}
\makesavenoteenv{longtable}
\usepackage{graphicx}
\makeatletter
\def\maxwidth{\ifdim\Gin@nat@width>\linewidth\linewidth\else\Gin@nat@width\fi}
\def\maxheight{\ifdim\Gin@nat@height>\textheight\textheight\else\Gin@nat@height\fi}
\makeatother
% Scale images if necessary, so that they will not overflow the page
% margins by default, and it is still possible to overwrite the defaults
% using explicit options in \includegraphics[width, height, ...]{}
\setkeys{Gin}{width=\maxwidth,height=\maxheight,keepaspectratio}
% Set default figure placement to htbp
\makeatletter
\def\fps@figure{htbp}
\makeatother
\setlength{\emergencystretch}{3em} % prevent overfull lines
\providecommand{\tightlist}{%
  \setlength{\itemsep}{0pt}\setlength{\parskip}{0pt}}
\setcounter{secnumdepth}{-\maxdimen} % remove section numbering
\ifLuaTeX
  \usepackage{selnolig}  % disable illegal ligatures
\fi
\usepackage{bookmark}
\IfFileExists{xurl.sty}{\usepackage{xurl}}{} % add URL line breaks if available
\urlstyle{same}
\hypersetup{
  pdftitle={RStudio Server on SPCS},
  hidelinks,
  pdfcreator={LaTeX via pandoc}}

\title{RStudio Server on SPCS}
\author{}
\date{\vspace{-2.5em}}

\begin{document}
\maketitle

Deploy \textbf{RStudio Server} as a custom \textbf{SPCS service} with
\textbf{skiPatrol} and \textbf{skiLift} pre-installed. This vignette
is the \textbf{technical reference} for the deploy kit shipped in this
package:

\begin{Shaded}
\begin{Highlighting}[]
\NormalTok{kit }\OtherTok{\textless{}{-}} \FunctionTok{system.file}\NormalTok{(}\StringTok{"rstudio{-}spcs"}\NormalTok{, }\AttributeTok{package =} \StringTok{"skiPatrol"}\NormalTok{)}
\FunctionTok{list.files}\NormalTok{(kit)}
\end{Highlighting}
\end{Shaded}

The Hitchhiker's Guide chapter \textbf{RStudio Server on SPCS} walks
through the same flow at a higher level.

\subsection{Architecture}\label{architecture}

\begin{verbatim}
Browser ──HTTPS ingress──► SPCS service (RStudio :8787)
                              ├── Native R session
                              ├── skiLift (SQL API / DBI)
                              ├── skiPatrol + reticulate
                              ├── Miniconda snowflake_ml (Snowpark Python)
                              └── /snowflake/session/token + SNOWFLAKE_HOST
\end{verbatim}

Unlike Workspace Notebooks, there is \textbf{no Python kernel} and no
\texttt{\%\%R} magic. R is the primary runtime.

\subsection{Deploy kit overview}\label{deploy-kit-overview}

\begin{longtable}[]{@{}
  >{\raggedright\arraybackslash}p{(\columnwidth - 2\tabcolsep) * \real{0.2727}}
  >{\raggedright\arraybackslash}p{(\columnwidth - 2\tabcolsep) * \real{0.7273}}@{}}
\toprule\noalign{}
\begin{minipage}[b]{\linewidth}\raggedright
Step
\end{minipage} & \begin{minipage}[b]{\linewidth}\raggedright
Script / file
\end{minipage} \\
\midrule\noalign{}
\endhead
\bottomrule\noalign{}
\endlastfoot
Configure & \texttt{config.example.env} → \texttt{config.env} \\
Provision stages / pool & \texttt{provision.sh} (from
\texttt{provision.sql.template}) \\
Build EAI (Image Builder) & \texttt{provision.sh\ -\/-eai} \\
RStudio password &
\texttt{RSTUDIO\_PASSWORD=\textquotesingle{}…\textquotesingle{}\ ./create\_password\_secret.sh}
(Snowflake secret) \\
Build image & \texttt{build\_snowflake.sh} or
\texttt{build\_local.sh} \\
Deploy service & \texttt{deploy\_service.sh} +
\texttt{service-spec.template.yaml} \\
Validate & \texttt{smoke\_test.R}, \texttt{smoke\_test.py} \\
\end{longtable}

All paths are under \texttt{inst/rstudio-spcs/} in the skiPatrol
repository.

\subsection{skiLift in RStudio}\label{skilift-in-rstudio}

\begin{Shaded}
\begin{Highlighting}[]
\FunctionTok{library}\NormalTok{(DBI)}
\FunctionTok{library}\NormalTok{(skiLift)}

\NormalTok{con }\OtherTok{\textless{}{-}} \FunctionTok{dbConnect}\NormalTok{(}\FunctionTok{Snowflake}\NormalTok{())}
\FunctionTok{dbGetQuery}\NormalTok{(con, }\StringTok{"SELECT CURRENT\_USER(), CURRENT\_WAREHOUSE()"}\NormalTok{)}
\end{Highlighting}
\end{Shaded}

Set \texttt{SNOWFLAKE\_WAREHOUSE}, \texttt{SNOWFLAKE\_DATABASE},
\texttt{SNOWFLAKE\_SCHEMA}, and \texttt{SNOWFLAKE\_ROLE} in the service
spec. See
\textbf{\texttt{vignette("spcs-custom-services",\ package\ =\ "skiLift")}}
for REST API limits (\texttt{USE\ WAREHOUSE},
\texttt{GET}/\texttt{PUT}).

\subsection{\texorpdfstring{skiPatrol:
\texttt{sfr\_connect\_spcs()}}{skiPatrol: sfr\_connect\_spcs()}}\label{skipatrol-sfr_connect_spcs}

\textbf{\texttt{sfr\_connect()}} auto-detects \textbf{Workspace
Notebook} sessions only. Custom SPCS services need explicit host + OAuth
token for Snowpark Python:

\begin{Shaded}
\begin{Highlighting}[]
\FunctionTok{source}\NormalTok{(}\StringTok{"\textasciitilde{}/spcs\_helpers.R"}\NormalTok{)   }\CommentTok{\# copied into image from the kit}
\NormalTok{conn }\OtherTok{\textless{}{-}} \FunctionTok{sfr\_connect\_spcs}\NormalTok{()}
\end{Highlighting}
\end{Shaded}

The helper:

\begin{enumerate}
\def\labelenumi{\arabic{enumi}.}
\tightlist
\item
  Points \textbf{reticulate} at Miniconda Python
  (\texttt{RETICULATE\_PYTHON}).
\item
  Sets \texttt{LD\_LIBRARY\_PATH} for Conda OpenSSL (avoids
  \texttt{OPENSSL\_3.3.0} mismatch).
\item
  Calls \texttt{skiPatrol}'s Python bridge with
  \texttt{authenticator\ =\ "oauth"}.
\end{enumerate}

\subsubsection{reticulate notes}\label{reticulate-notes}

\begin{itemize}
\tightlist
\item
  Use \textbf{\texttt{reticulate::use\_python(full\_path)}} --- not
  \texttt{use\_condaenv()}.
\item
  The reference image installs \textbf{Miniconda} env
  \texttt{snowflake\_ml} at
  \texttt{/opt/conda/envs/snowflake\_ml/bin/python}.
\item
  Run \texttt{conda\ install\ openssl} in that env if you extend the
  Dockerfile.
\end{itemize}

\subsubsection{\texorpdfstring{Snowpark \texttt{collect()} vs
pandas}{Snowpark collect() vs pandas}}\label{snowpark-collect-vs-pandas}

Snowpark \textbf{\texttt{collect()}} returns a list of \texttt{Row}
objects, not a pandas DataFrame:

\begin{Shaded}
\begin{Highlighting}[]
\NormalTok{rows }\OperatorTok{=}\NormalTok{ session.sql(}\StringTok{"SELECT 1 AS x"}\NormalTok{).collect()}
\NormalTok{rows[}\DecValTok{0}\NormalTok{].as\_dict()}
\end{Highlighting}
\end{Shaded}

For R via reticulate, prefer \textbf{\texttt{to\_pandas()}} or
\texttt{py\_to\_r()} on the result.

\subsection{Smoke tests}\label{smoke-tests}

After deploy, log in via ingress SSO, then RStudio (\texttt{rstudio} /
password from the Snowflake secret):

\begin{Shaded}
\begin{Highlighting}[]
\FunctionTok{source}\NormalTok{(}\StringTok{"\textasciitilde{}/smoke\_test.R"}\NormalTok{)}
\end{Highlighting}
\end{Shaded}

\begin{Shaded}
\begin{Highlighting}[]
\ExtensionTok{python3}\NormalTok{ \textasciitilde{}/smoke\_test.py}
\end{Highlighting}
\end{Shaded}

\subsection{Image build: tarballs}\label{image-build-tarballs}

Compiling R packages inside Image Builder is slow. The kit's
\texttt{prepare\_build\_ctx.sh} builds a \textbf{skiPatrol} tarball
locally and copies it into the flat build context. \textbf{skiLift}
is included when available; otherwise \texttt{install\_packages.R}
installs from GitHub (requires build EAI).

For production, pin tarball URLs or commit \texttt{.tar.gz} files under
\texttt{pkg/} in the kit directory.

\subsection{Ingress URL}\label{ingress-url}

The public endpoint hostname \textbf{changes} when you
\texttt{DROP\ SERVICE} and \texttt{CREATE\ SERVICE} again. Always run:

\begin{Shaded}
\begin{Highlighting}[]
\NormalTok{SHOW ENDPOINTS }\KeywordTok{IN}\NormalTok{ SERVICE MY\_DB.MY\_SCHEMA.RSTUDIO\_SVC;}
\end{Highlighting}
\end{Shaded}

\subsection{Persistence}\label{persistence}

Mount internal stages for workspace files and user-installed R packages:

\begin{Shaded}
\begin{Highlighting}[]
\FunctionTok{volumes}\KeywordTok{:}
\KeywordTok{{-}}\AttributeTok{ }\FunctionTok{name}\KeywordTok{:}\AttributeTok{ workspace}
\AttributeTok{  }\FunctionTok{source}\KeywordTok{:}\AttributeTok{ }\StringTok{"@MY\_DB.MY\_SCHEMA.VOLUMES/rstudio\_workspace"}
\KeywordTok{{-}}\AttributeTok{ }\FunctionTok{name}\KeywordTok{:}\AttributeTok{ r{-}libs}
\AttributeTok{  }\FunctionTok{source}\KeywordTok{:}\AttributeTok{ }\StringTok{"@MY\_DB.MY\_SCHEMA.VOLUMES/rstudio\_r\_libs"}
\end{Highlighting}
\end{Shaded}

\subsection{When not to use this
pattern}\label{when-not-to-use-this-pattern}

\begin{longtable}[]{@{}ll@{}}
\toprule\noalign{}
Alternative & Why \\
\midrule\noalign{}
\endhead
\bottomrule\noalign{}
\endlastfoot
Workspace + \texttt{setup\_notebook()} & No Docker; notebook-first \\
Posit Workbench Native App & Managed multi-user IDE \\
Snowflake Remote Dev & VS Code / Cursor; no RStudio \\
\end{longtable}

\subsection{Related reading}\label{related-reading}

\begin{itemize}
\tightlist
\item
  \textbf{\texttt{vignette("spcs-custom-services",\ package\ =\ "skiLift")}}
  --- DBI auth and REST API constraints
\item
  \textbf{\texttt{vignette("workspace-notebooks")}} --- \texttt{\%\%R} /
  bootstrap path
\item
  Kit README:
  \texttt{system.file("rstudio-spcs/README.md",\ package\ =\ "skiPatrol")}
\end{itemize}

\end{document}
